\documentclass{nfr}

% This LaTeX template is based on the official Word template found here: https://www.forskningsradet.no/siteassets/utlysninger/vedlegg-utlysninger/template-for-project-description-researcher-project-for-experienced-scientists-and-early-career-scientists.docx
% See the README.md file for more information

% Here's an example bibliography using the numeric (Nature) style and two
% columns for the reference list.

% \usepackage{multicol}
% \renewcommand{\bibpreamble}{\begin{multicols*}{2}}
% \renewcommand{\bibpostamble}{\end{multicols*}}

\usepackage{comment}

% \usepackage{biblatex}
% \addbibresource{references.bib}

% \citestyle{nature}
% \citestyle{}

% \bibliographystyle{bst/veryshortnum}
\bibliographystyle{apalike}

\title{Teaching Statement}
\author{Zhongli Wang}

\begin{document}

% Project description template
% Researcher Project for Experienced Scientists/Early Career Scientists

% This template consists of
% 
% \textbf{A:} Overview of the chapters and sections to be used as the
% structure of the project description
% 
% \textbf{B:} General guidance
% 
% \textbf{C:} Guidance on the content in each chapter and section
% 
% \textbf{A: Chapters and sections to be used}

% I used pandoc to translate the word template to tex. The formatting of the
% appendices are, therefore, not great, but the contents are there.

\maketitle

My teaching philosophy centers on clarity, inclusion, and interaction. I view teaching as a process of mutual learning in which both students and instructors advance through curiosity, dialogue, and structured reasoning. My goal is to help students build analytical confidence by connecting abstract economic concepts to real-world applications through clear explanations and interactive exercises. By demonstrating how data and code can make complex economic mechanisms tangible, I aim to cultivate students’ intuition and technical competence simultaneously. These tools not only enhance comprehension but also prepare students for the modern, data-rich environment of economic research and policy analysis.

\begin{center}
\small{
\begin{tabular}{||c c c c||} 
 \hline
 Course Name & Institute & Lecturer & Semester \\ [0.5ex] 
 \hline\hline
 Applied Empirical Methods for Business Administration & UZH & Prof. Karl Schmedders & 2019 Fall \\ 
 \hline
 Applied Business Modelling and Analytics, Part 1 & UZH & Dr. Maik Meusel & 2020 Spring \\
 \hline
\end{tabular}}
\end{center}

During my master’s studies at the University of Zurich, I worked part-time for two years as a teaching assistant at the Chair of Quantitative Business Administration (QBA). I contributed to the SYLVA project, an automated scoring and feedback system designed to enhance the efficiency and quality of instruction in Economics, Finance, and Business Administration. My responsibilities included developing and refining Mathematica code, testing algorithmic accuracy, and assisting lecturers in improving course materials. I also supported students by answering questions about the system’s logic, the interpretation of results, and lecture-related exercises. This experience strengthened my understanding of how digital tools can complement human instruction, promote fairness in evaluation, and make learning more interactive and engaging.\\

I believe in creating an inclusive and supportive learning environment, where students feel encouraged to ask questions freely and where diverse perspectives serve as catalysts for discovery. Beyond the classroom, I am eager to collaborate with colleagues on developing shared teaching materials, organizing workshops on computational methods, and mentoring students interested in research. I view teaching not merely as knowledge transfer but as a collaborative endeavor that fosters intellectual curiosity, critical thinking, and lifelong learning for both students and educators. In October 2025, together with my colleague Dr.~Hankui Wang, I delivered a public lecture on “AI Agents in Science” at the Wissenswerkstadt Bielefeld, a science museum dedicated to public engagement with research and technology. Presenting in German, I introduced the concept of agentic workflows and large language models as emerging tools for economics research and education. Engaging with the audience’s questions in real time reinforced the importance of clarity, patience, and flexibility when communicating complex ideas to people with diverse technical backgrounds.\\

Looking ahead, I am prepared to teach a wide range of undergraduate and graduate courses in economics, including microeconomics, industrial organization, and applied econometrics. I have a particular interest in designing courses that integrate computational tools and simulation-based methods into economic analysis. In future teaching, I plan to combine theoretical insight with hands-on practice by incorporating reproducible research tools, data-driven case studies, and interdisciplinary examples. My goal is to make abstract economic concepts both tangible and relevant, helping students develop not only technical competence but also the curiosity and confidence to apply economic reasoning to real-world and policy challenges.

\begin{comment}
\section*{Current Research}

% Here's a citation to \cite{texbook}. Here are citations to \cite{texbook,knuth:1984,lesk:1977}, and to \cite{lesk:1977,texbook,latex2e,latex:companion}.

My doctoral dissertation, ``\textit{Strategic Behavior and Economic Dynamics in ICT Markets: An Agent-based Framework}", investigates how firms in the information and communications technology (ICT) industry adapt under constraints such as data regulations, capacity limits, and consumer network effects. The work builds a bridge between state-space and agent-based models, showing how techniques from dynamic latent-variable systems can be applied to simulate and estimate complex economic environments.\footnote{All three chapters are single-authored papers.}\\

My job market paper, \textit{``Strategic Firm Capacity Selection under Network Effects: A Computational Model of Capacity Investment Dynamics in the ICT Industry''}, investigates how digital firms determine optimal service capacity under the joint influence of consumer congestion and network effects---an interaction that has received limited attention in the existing literature. While traditional industrial organization models emphasize innovation or price competition, many digital firms face a different constraint: their profitability depends critically on infrastructure investment and the ability to maintain service quality. To address this challenge, I combine game-theoretic reasoning with agent-based simulations. The game-theoretic framework captures the multi-stage nature of firms’ decisions on pricing and capacity investment, while the agent-based model represents the market dynamics that emerge when heterogeneous consumers interact through a small-world network and share information within user communities. The results show that strong network effects tend to amplify competitive advantages by concentrating demand, whereas congestion generates counteracting pressures that reduce pricing flexibility and erode profits. Together, these forces explain why digital firms must balance capacity expansion and user experience to sustain long-term profitability. This study is among the first to systematically analyze the joint role of network and congestion effects in the ICT industry, offering new insights into pricing, investment, and regulatory strategies for digital service markets.\\

Building on the similar market context, my previous paper, \textit{``Heterogeneous Consumers, Firm Competition, and the Data Value Chain: An Agent-Based Approach''}, investigates how differences in data-governance intensity influence competition in software markets. The challenge arises because digital firms depend on user data to enhance product quality and pricing decisions, while local governments impose varying levels of data regulation to protect privacy and maintain market fairness. 
To analyze this tension, I develop a hybrid game-theoretic and agent-based model in which governments determine the stringency of data governance, and firms subsequently compete on investment and pricing decisions. 
Consumers are heterogeneous in their sensitivity to product quality and their attitudes toward data governance, which shape their purchasing behavior. 
The results show that stricter data governance limits firms’ ability to utilize consumer data effectively and reduces competitive differentiation, which leads to higher prices and lower consumer welfare. In contrast, moderate levels of regulation can maintain competition and promote more balanced market outcomes. The analysis suggests that regional governments would benefit from greater coordination in setting data policies and from incorporating consumer preferences into policy design, thereby supporting both market efficiency and innovation. Overall, these findings offer a structured perspective on how policymakers can balance privacy protection with incentives for innovation in digital markets.\\

Building on the insights from the first two studies, my third paper ``\textit{Parameter Estimation Methods in an Agent-based Computational Model: Insights on ICT Industrial Dynamics and Policies}" focuses on empirical validation. 
Using financial and operational data from two leading U.S. telecommunication firms, AT\&T and T-Mobile, I estimate the model parameters with the Simulated Method of Moments and neural network–based surrogate modeling (both within-sample and out-of-sample check). 
The main challenge lies in linking complex agent-based models to real-world data, as repeated simulations are computationally intensive and traditional estimation methods often fall short. 
To address this, I employed the Paderborn High-Performance Computing (HPC) clusters and the Python package \texttt{ray} to parallelize the simulations and collect 40{,}000 samples for parameter search. 
By comparing these two approaches, the study provides practical evidence on how to calibrate computational models efficiently while maintaining empirical accuracy. This work contributes to the broader methodological discussion on integrating simulation-based models with data-driven analysis to better understand market behavior and policy outcomes.



\section*{Publications and Emerging Projects}

In addition to my market-oriented research, I have one peer-reviewed publication and one current working paper that contribute to the broader use of computational methods in economics.\footnote{Both papers are coauthored with other researchers.} \\

The paper \textit{``Polytope Fraud Theory''} extends the classical Fraud Triangle and Fraud Diamond frameworks by introducing a multidimensional structure that captures ten abnormal financial practice alarms indicative of potential corporate fraud \citep{Zhao2025}. Using data on activist short sellers and supervised machine learning techniques, the study identifies financial patterns associated with fraudulent reporting and demonstrates how algorithmic approaches can support market transparency and investor protection. 
The paper also demonstrates the potential for automating advanced financial statement analysis, providing tools that can enhance both auditing efficiency and investment decision-making. In addition, it introduces the \textit{Unified Investor Protection Framework}, which systematically categorizes investor-protection theories from the macro- to the micro-level. This framework offers practical guidance for identifying firm-level financial risks and strengthens the connection between analytical modeling and real-world investor decision processes.\\

Another working paper, \textit{``Agentic Workflows for Economic Research: Design and Implementation''}, develops a framework that employs large-language-model (LLM) agents to support economists throughout the research process—from data collection and coding to simulation design and interpretation \citep{Dawid2025}. The framework formalizes how autonomous AI agents can collaborate within multi-stage research pipelines while ensuring reproducibility, version control, and transparent documentation. Using examples implemented in Microsoft AutoGen, we demonstrate how such workflows can accelerate tasks like research-question ideation, literature review, and model specification, thereby improving research efficiency without compromising academic rigor. Importantly, the framework positions human researchers at the core of the process: AI agents act as supportive tools that augment, rather than replace, human judgment and creativity. This line of work extends my methodological interests in AI-enhanced modeling, bridging computational economics with artificial-intelligence research and tool development.

\section*{Future Research}

\subsection*{Potential Research Direction: The Observability of Firm Heterogeneity in Price-Directed Consumer Search}

% Recent models of price-directed search \citep{obradovits2023price} highlight how consumers allocate their attention and search effort across differentiated firms when facing information frictions. However, these frameworks often assume that firm heterogeneity---in cost structures, quality, or service intensity---is either fully observable to consumers or symmetrically distributed across sellers. In many real markets, such as digital platforms, online retailing, or app ecosystems, this heterogeneity is only partially observable: consumers infer firm characteristics indirectly from prices, reputation indicators, or platform signals.

% Understanding how the observability of firm heterogeneity shapes consumer search behavior and market efficiency is crucial for both theory and policy. It determines the extent of equilibrium price dispersion, the efficiency of market matching, and firms’ incentives to disclose information. The problem is empirically feasible because modern datasets (e.g., online pricing panels, platform APIs, or laboratory experiments) provide micro-level evidence on search behavior under varying degrees of information transparency. Computational approaches, including agent-based simulations, further enable the exploration of equilibrium dynamics in complex environments.

Recent models of price-directed search \citep{obradovits2023price} examine how consumers allocate attention and search effort across differentiated firms under information frictions, yet they typically assume that firm heterogeneity---in costs, quality, or service intensity---is fully observable or symmetrically distributed. In real markets such as digital platforms or online retailing, this heterogeneity is only partially observable, as consumers infer firm characteristics indirectly from prices, ratings, or platform signals. Understanding how such observability affects consumer search and market efficiency is crucial, as it shapes equilibrium price dispersion, allocative outcomes, and firms’ incentives to disclose information. The question is empirically feasible given modern micro-level data on online search and pricing, while computational approaches, including agent-based simulations, allow systematic analysis of equilibrium dynamics under varying information transparency. \\

% This research direction contributes to Industrial Organization by linking information structure, firm heterogeneity, and price competition under bounded rationality. It combines theoretical modeling with simulation and empirical validation, making it both conceptually rich and methodologically tractable. 

The following hypotheses will be tested:

\begin{enumerate}%[label=\textbf{H\arabic*.}, leftmargin=2em]
    \item \textbf{Partial observability increases price dispersion.} When firm heterogeneity (e.g., in cost or quality) is imperfectly observable, consumers engage in less directed search, resulting in a wider equilibrium price distribution than under full observability.
    
    \item \textbf{High observability improves allocative efficiency.} Markets with more transparent signals of firm heterogeneity (such as ratings or quality certifications) achieve a closer match between consumers and preferred firms, raising consumer surplus and reducing welfare loss.
    
    \item \textbf{Strategic transparency depends on firm type.} Low-cost or high-quality firms have stronger incentives to increase the observability of their characteristics (e.g., through signaling or advertising) than high-cost or low-quality competitors, leading to asymmetric information equilibria.
\end{enumerate}

\subsection*{Methodological Plan}

In my approach, I plan to combine analytical and computational methods. Analytical derivations will be used to establish the basic propositions and benchmark outcomes of the model, while numerical techniques and computational modeling will be employed to solve for Nash equilibria and explore market dynamics under varying degrees of information transparency. For empirical validation, I have access to the Refinitiv database through an application via the Southern University of Science and Technology (SUSTech) and may also rely on additional database resources at my future host institute. My doctoral research was supported by the Marie Skłodowska-Curie Actions and the Bielefeld Nachwuchsfonds, and it benefited from extensive computational work conducted on Google Colab (self-funded) and the Paderborn High Performance Computing cluster (supported by Bielefeld University). I hope that the new institute could provide adequate support for research infrastructure, particularly computational resources, to further advance this line of study.



% \subsection*{Computational \& ABM Dynamics}

% Specifically, I plan to:

% •	Develop scalable frameworks for data-driven research of economic models that can inform competition and innovation policy.

% •	Study how AI and automation transform firm strategy, market structure, and industrial dynamics.

% My long-term objective is to contribute to the economics literature, that is both analytically rigorous and computationally adaptive—one that captures the rich, evolving behavior of firms and consumers in real markets. I remain deeply motivated to pursue research that connects theoretical innovation with practical relevance, and I look forward to collaborating in an environment that values interdisciplinary approaches and empirical credibility.
% \citep{arve2025decentralized}.


% \bibliography{NFR_Template_Statements/references}
\bibliography{NFR_Template_Statements/references1}

\newpage

\appendix
\end{comment}

\end{document}
